\documentclass{article}
\usepackage{amsmath}
\usepackage{amssymb}
\usepackage{amsthm}
\usepackage{MnSymbol}
\usepackage{xcolor}
\usepackage{parskip}
\usepackage{tikz}
\usepackage{bbm}
\usepackage{hyperref}
\usetikzlibrary{trees}
\usetikzlibrary{graphs}


\theoremstyle{definition}
\newtheorem{definition}{Definition}
\newtheorem{theorem}{Theorem}


\newcommand{\abs}[1]{\left| #1 \right|}


\title{lecture}
\date{\today}
\begin{document}
\maketitle

\section{approximation resistance}
if an algorithm is \emph{approximation resistant}, then it has an upper bound in how well a polynomial time problem can approximate the best solution. beating the upper bound would prove \(\mathbb P=\mathbb{NP}\)



imagine a ranking. you can query an oracle for rankings of two people and it tells you which one is greater. trivial random ordering gives 0.5 worst case.

use SDP

make a graph where if \(a<b\), then there is a directed edge w weight 1 from \(a\) to \(b\), and a directed edge with weight \(-1\) from \(b\) to \(a\). max cut gives a better than optimal configuration for each half. 

if the right side of maxcut has more +'s going in, then it becomes the right half of the ordering. recurse. 









\end{document}